\section{Concurrency level vs.~observed parallelisam}
\label{section:node-level:concurrency}

\begin{hypothesis}
 The concurrency level is, in principle, high, but fails to materialise.
 The observes parallelism hence is a result of other flaws, but not of a lack of actual parallelisation in the code.
\end{hypothesis}


Jetzt haben wir ja ganz bewusst weniger Tasks erzeugt. 
Somit ist the Evaluierung mit Kontext zu versehen
Task Begriff gilt hier auch fuer loop iterations.

\begin{hypothesis}
 There are too many critical sections that cannot be entered by more than one thread, such that the threads tend to synchronise each other.
\end{hypothesis}


\paragraph{Data collection}

\begin{itemize}
  \item The effective CPU time (in the user code) is low compared to the actual execution time.
  \item There is a huge spin time.
  \item We see a massive lock contention.
\end{itemize}

\noindent

\begin{instructions}{\vtune}
In \vtune, too many synchronisation statements materialise in a high spin time as compared to the effective time.  
This leads to a low effective CPU utilisation (below is an example with eight cores, where we effectively only use four of them).
\end{instructions}


\subsection{Evaluation}

\begin{enumerate}
  \item[$\boxtimes$] In line with all other metrics, we consider a code phase or loop reasonably balanced, if it is able to exploit around 80\% of the available compute resources at any point in time.
  \item[$\boxtimes$] If we underrun this threshold, the code is not super balanced.
\end{enumerate}

\paragraph{Evaluation}

If the spin time or overhead time in the area of interest exceeds 20\% of the total execution time, i.e.~if the time we actually spend in our code underruns 80\% of the total time, we have a severe synchronisation flaw.


\paragraph{Optimisation and follow-up steps}



\paragraph*{How it works}

\begin{itemize}[noitemsep]
  \item[$\boxplus$] Validate that the administrative (non-science) fraction of the code scales with the scientific core calculations.
\end{itemize}


\paragraph*{How it does not work}

\begin{itemize}[noitemsep]
  \item[$\boxminus$] The administrative overhead will always increase relatively, so an increase in itself is not a flaw. An extraordinary increase coupled with limited scalability however is.
\end{itemize}
